\documentclass[12pt]{article}
\usepackage[utf8]{inputenc}
\usepackage[T1]{fontenc}
\usepackage{hyperref}
\title{Zialo Language Resources for K-12 Curriculum Development}
\author{Compiled by Claude Sonnet 5.5 (claude-sonnet-5-5) for Andr\'as Kornai}
\date{}
\begin{document}
\maketitle

\section{Introduction}
Zialo (ISO 639-3: zil; Glottocode: zial1235) is a Southwestern Mande language of southeastern Guinea, mainly in Macenta prefecture. Estimates are about 25,000 speakers. It was long treated as a dialect of Loma and was identified as a distinct language by Kirill Babaev's fieldwork (2010). Sources report no established writing system, and French is the language of official use.

\section{Primary Educational Resources}
No primers, textbooks or school programmes in Zialo were found in search results.

\section{Linguistic and Reference Materials}
\begin{itemize}
\item Wikipedia, Zialo language: \url{https://en.wikipedia.org/wiki/Zialo_language}
\item Wikipedia, Zialo people: \url{https://en.wikipedia.org/wiki/Zialo_people}
\item Endangered Languages Project, with information from Babaev (2010), \emph{Zialo: the newly-discovered Mande language of Guinea} (Lincom Europa), described as the first grammar of the language: \url{https://endangeredlanguages.com/lang/4573}
\item Endangered Languages Project entry for Zialo: \url{https://endangeredlanguages.com/node/92516}
\item LINGUIST List announcement of the Babaev book: \url{https://linguistlist.org/issues/21/4220/}
\end{itemize}

\section{Religious and Community Resources}
None found in search results.

\section{Digital and Audio Resources}
None found in search results.

\section{Curriculum Development Context}
Zialo has a documented grammar but no known orthography or educational materials. Any mother-tongue curriculum would begin with orthography design and community consultation, using French as the school language of transition.

\end{document}